\documentclass[a4paper,12pt]{letter}
% POZOR: samo utf8. V starih datotekah je bil za njim se \usepackage[cp1250]{inputenc},
% kar je prevzelo prednost in lahko pokvari sumnike.
\usepackage[utf8]{inputenc}
\usepackage[T1]{fontenc}
\usepackage[english,slovene]{babel}
\usepackage[tikz]{bclogo}
\usepackage{qrcode,marvosym}
\usepackage{pgf,tikz,pgfplots}
\pgfplotsset{compat=1.14}
\usepackage{mathrsfs}
\usetikzlibrary{arrows}
\usepackage{graphicx}
\newcommand{\ds}{\displaystyle}
\usepackage{fancybox}
\usepackage{amsfonts}
\usepackage{amsmath}
\usepackage{wallpaper}
\usepackage{hyperref}
\usepackage{array}
\newcolumntype{M}[1]{>{\centering\arraybackslash}m{#1}}
\newcolumntype{N}{@{}m{0pt}@{}}
\usepackage{fancyhdr,enumerate}
\newcommand{\ora}{\overrightarrow}
\renewcommand{\headrulewidth}{3pt}
\renewcommand{\headrule}{\hbox to\headwidth{%
  \color{gray}\leaders\hrule height \headrulewidth\hfill}}
\renewcommand{\footrulewidth}{2pt}
\renewcommand{\footrule}{\hbox to\headwidth{%
  \color{brown}\leaders\hrule height \footrulewidth\hfill}}
\usepackage{gensymb}
\usepackage{amssymb}
  \textwidth 20 true cm
  \oddsidemargin -2 true cm
  \evensidemargin -2 true cm
  \topmargin -3.5 true cm
  \textheight 27.5 true cm
  \linespread{1.2}
\renewcommand{\headrulewidth}{0.3pt}
\renewcommand{\footrulewidth}{1.9pt}
\definecolor{qqwuqq}{rgb}{0.,0.39215686274509803,0.}
\definecolor{qqqqff}{rgb}{0,0,1}
\definecolor{ffqqqq}{rgb}{1.,0.,0.}
\definecolor{zzuxux}{rgb}{0,0,0}
\usepackage{mdframed}
\usepackage{multicol}
\definecolor{bgblue}{RGB}{0,0,253}
\usepackage{eurosym}
%%%%%%%%%%%%%%%%%%%%%%%%%%%%%%%%%%%%%%%%%%%%%%%%%%%%%%%%
\newcounter{tocke}
\setcounter{tocke}{0}
\usepackage{tikz-3dplot}
%%%%%%%%%%%%%%%%%%%%%%%%%%%%%%%%%%%%%%%%%%%%%%%%%%%%%%%%%%
\mdfdefinestyle{theoremstyle}{%
linecolor=brown!40,linewidth=1pt,%
frametitlerule=true,%
frametitlebackgroundcolor=brown!50,
innertopmargin=\topskip,
}
\mdfdefinestyle{theoremstyle1}{%
linecolor=brown,middlelinewidth=1pt,%
frametitlerule=true,%
apptotikzsetting={\tikzset{mdfframetitlebackground/.append style={%
shade,left color=black!40, right color=gray!10},
mdfbackground/.append style={
top color=white!100!white,
bottom color=black!5
},
}},
frametitlerulecolor=gray,
frametitlerulewidth=2pt,
innertopmargin=\topskip,
innerbottommargin=\topskip,
}
\mdtheorem[style=theoremstyle1]{naloga}{Naloga}
\mdtheorem[style=theoremstyle]{resitev}{Re\v{s}itev}
\mdtheorem[style=theoremstyle1]{kriterij}{Kriterij ocenjevanja}
%%%%%%%%%%%%%%%%%%%%%%%%%%%%%%%%%%%%%%%%%%%%%%%%%%%%%%%%%%
\usepackage[table]{xcolor}
\usepackage[printwatermark]{xwatermark}
\newwatermark[allpages,color=brown!2,angle=45,scale=5,
xpos=-5,ypos=0]{matej.info}
%%%%%%%%%%%%%%%%%%%%%%%%%%%%%%%%%%%%%%%%%%%%%%%%%%%%%%%%%%

\begin{document}
\pagestyle{fancy}
\renewcommand\headrulewidth{0pt}
\lhead{}\chead{}\rhead{}
\rfoot{\vspace*{-2\baselineskip}}
\renewcommand{\vec}{\overrightarrow}
\everymath{\displaystyle}

\begin{minipage}{20cm}
\begin{bclogo}[couleur=brown!40,
  arrondi=0,ombre=true, couleurOmbre=gray!50,
logo =\bcplume,
  noborder=true]
{\textrm{{\Large{PREVERJANJE 1.P.2} \Huge{}:
\large{Eksponentna funkcija. Logaritemska funkcija} \hfill $G-3$}}}
\end{bclogo}
\end{minipage}
\begin{minipage}{20cm}
\begin{bclogo}[couleur=brown!40,ombre=true, couleurOmbre=gray!50,
  arrondi=0,
logo =\bcplume,
  noborder=true]
{\textsc{{\large{Ime in priimek:}
\vspace{0.2cm} \rule{15cm}{0.2mm}\hfill }}}
\end{bclogo}
\end{minipage}


%%%%%%%%%%%%%%%%%%%%%%%%%%%%%%%%%%%%%%%%%%%%%%%%%
\def\tockeA{4+4+4}
\begin{naloga}[\hfill $\tockeA$] \addtocounter{tocke}{\numexpr\tockeA\relax}
Re\v{s}i ena\v{c}bo:

a.) $\log_2(x+2)+\log_2(x-2)=5$

b.) $5^{\frac{x}{2}}\cdot\sqrt{125}=25^{\,x-1}$

c.) $2^{x-1}=7^{x}$ \quad (re\v{s}itev zapi\v{s}i v obliki logaritma $\log_a b$)
\end{naloga}\newpage

%%%%%%%%%%%%%%%%%%%%%%%%%%%%%%%%%%%%%%%%%%%%%%%%%
\def\tockeB{2+1+1+3}
\begin{naloga}[\hfill $\tockeB$] \addtocounter{tocke}{\numexpr\tockeB\relax}
Podana je funkcija $f(x)=3^{\,x+1}-9$.

a.) Dolo\v{c}i ni\v{c}lo, za\v{c}etno vrednost in zalogo vrednosti funkcije $f$.

b.) Zapi\v{s}i inverzno funkcijo $f^{-1}$ in njeno definicijsko obmo\v{c}je.
\end{naloga}\newpage

%%%%%%%%%%%%%%%%%%%%%%%%%%%%%%%%%%%%%%%%%%%%%%%%%
\def\tockeC{4}
\begin{naloga}[\hfill $\tockeC$] \addtocounter{tocke}{\numexpr\tockeC\relax}
Izra\v{c}unaj brez kalkulatorja:
$$\log_a a^5+\log_a \sqrt{a^3}-\log_2 7a\cdot\log_{7a} 16$$
\end{naloga}\newpage

%%%%%%%%%%%%%%%%%%%%%%%%%%%%%%%%%%%%%%%%%%%%%%%%%
\def\tockeD{4}
\begin{naloga}[\hfill $\tockeD$] \addtocounter{tocke}{\numexpr\tockeD\relax}
Na sliki je graf funkcije $f(x)=\log_b(x+a)+c$ z navpi\v{c}no asimptoto $x=-2$.
Dolo\v{c}i konstante $a$, $b$ in $c$ ter izra\v{c}unaj $\dfrac{f(25)+1}{f(7)+1}$.

\begin{center}
\begin{tikzpicture}[line cap=round,line join=round,>=triangle 45,x=1.0cm,y=1.0cm]
\begin{axis}[
x=1.0cm,y=1.0cm,
axis lines=middle,
ymajorgrids=true,
xmajorgrids=true,
xmin=-3.5,xmax=8.5,
ymin=-3.5,ymax=3.5,
xtick={-3.0,-2.0,...,8.0},
ytick={-3.0,-2.0,...,3.0},]
\clip(-3.5,-3.5) rectangle (8.5,3.5);
\draw[line width=1.pt,color=qqwuqq,smooth,samples=100,domain=-1.98:8.5] plot(\x,{ln((\x)+2)/ln(3)-1});
\draw [line width=1.pt,dash pattern=on 2pt off 2pt,color=ffqqqq] (-2,-3.5) -- (-2,3.5);
\begin{scriptsize}
\draw[color=qqwuqq] (7.7,2.9) node {$f$};
\end{scriptsize}
\end{axis}
\end{tikzpicture}
\end{center}
\end{naloga}

\newpage
\vfill
\begin{kriterij*}[\hfill \v{s}tevilo vseh to\v{c}k na preverjanju: $\arabic{tocke}$]
 $$
 \begin{array}{||c|c|c|c|c|c||>{}c|c|}
 \hline\hline
 \textup{ocena}& 1&2&3&4&5&\textrm{uspe\v{s}nost v } \%& \textrm{\bf OCENA}\\
  \hline\hline
 \% & [0,45) &[45,60)&[60,75) &[75,90) & [90,100]& \mbox{\phantom{ \huge{$\frac{5}{445}$ 15}}} & \\
 \hline
 \end{array} \hspace*{0.5cm}\vspace*{-0.3cm} \qrcode[hyperlink,height=0.8in]{http://www.matej.info}
$$
\vspace*{0.1cm}
\end{kriterij*}

\newpage
\section*{Kon\v{c}ne re\v{s}itve}

\begin{enumerate}
\item a.) $x=6$ \quad b.) $x=\dfrac{7}{3}$ \quad c.) $x=\log_{\frac{2}{7}} 2$
\item a.) ni\v{c}la $x=1$, za\v{c}etna vrednost $f(0)=-6$, zaloga vrednosti $(-9,\infty)$ \quad b.) $f^{-1}(x)=\log_3(x+9)-1$, $D_{f^{-1}}=(-9,\infty)$
\item $\dfrac{5}{2}$
\item $a=2$, $b=3$, $c=-1$; \quad $\dfrac{f(25)+1}{f(7)+1}=\dfrac{3}{2}$
\end{enumerate}
\end{document}
